% DRAFT generated by akm.run draft. Sentences below were CHOSEN by the outcomes;
% if a prediction failed, a TODO marks where your explanation must go.

\label{sec:results}
\paragraph{Detection (Experiment 1).} Table~\ref{tab:e1} gives the full matrix.
Keyless verification was blind to key-transformation faults: V0 detected \akmVzeroWrongKey{} and V1 \akmVoneWrongKey{} of \akmWrongKeyTrials{} wrong-key trials (F5, F5b, F5c), while V2$^{-}$, V2 and V3 detected \akmVtwoMinusWrongKey{}, \akmVtwoWrongKey{} and \akmVthreeWrongKey{} respectively.
Every verifier detected every structural fault (skips, duplicates, bit-flips, edited receipts and roots) in every trial.
On \akmFzeroTrials{} honest verifier-trials, no verifier raised a false alarm.
After retirement of $K_\text{old}$ (F7), V2$^{-}$ could not run while V2 still detected the fault; F7 was the only row separating the two (P6).
After a crash between receipt and flip (F6), every verifier reported the migration as incomplete with the exact set of flipped records.
All pre-registered Experiment~1 predictions held.

\paragraph{Sampling (RQ4).} 
At $\varepsilon=1\%$ and $k=100$, V1-S detected skipped records in \akmRqFoneb{}\% and bit-flipped records in \akmRqFthree{}\% of trials, against the textbook $1-(1-\varepsilon)^k=\akmRqTheory{}\%$.
On wrong-key faults it detected nothing at any $k$ or $\varepsilon$: across the sweep, V1-S directly challenged \akmRqFfiveChallenged{} records carrying a wrong key and accepted every one, exactly as Corollary~1.1 predicts.
No innocent record was ever flagged (\akmRqInnocent{} false flags).
Of the \akmRqChecked{} skip and bit-flip cells, \akmRqMissNinetyFive{} fell outside their 95\% interval (about \akmRqMissExpected{} are expected by chance, more because cells share challenge sequences), and \akmRqMissBonf{} fell outside the pre-registered Bonferroni-adjusted interval.
The single exception was F3 at $\varepsilon=5\%$, $k=200$, where the observed rate (99.9\%) fell just outside the Bonferroni-adjusted interval against a theoretical prediction of 99.997\%. At this near-saturated detection level, Wilson intervals are known to be fragile; critically, the exactness check above confirms zero missed detections and zero false flags throughout the entire sweep, indicating this is a boundary artifact of finite-sample estimation at an extreme probability, not a genuine gap in detection capability.

\paragraph{Cost (Experiment 2).} Tables~\ref{tab:e2time}--\ref{tab:p4} give the timings.
The ratio of V2 verification to proof-emitting migration (P4, threshold 0.10) was: per-record fsync, wall-clock: \akmPfourperrecordwallsmin{}--\akmPfourperrecordwallsmax{} (4 of 4 configurations below 0.10); per-record fsync, CPU time: \akmPfourperrecordcpusmin{}--\akmPfourperrecordcpusmax{} (4 of 4 configurations below 0.10); group commit, wall-clock: \akmPfourgroupOneZeroZeroZerowallsmin{}--\akmPfourgroupOneZeroZeroZerowallsmax{} (0 of 4 configurations below 0.10); group commit, CPU time: \akmPfourgroupOneZeroZeroZerocpusmin{}--\akmPfourgroupOneZeroZeroZerocpusmax{} (0 of 4 configurations below 0.10).
P4 therefore holds only where per-record fsync dominates migration wall-clock time; on CPU time or under group commit, verification is a substantial fraction of migration, as anticipated directly in the pre-registration.
Emitting the proof (receipts and Merkle bundle) added \akmOverheadRecMin{}--\akmOverheadRecMax{}\% to migration time under per-record fsync and \akmOverheadGrpMin{}--\akmOverheadGrpMax{}\% under group commit.

\paragraph{Storage.} Table~\ref{tab:storage} gives the per-record overhead.
The receipt log costs \akmReceiptBytes{}\,B per record for every rung, the commitment field only $\tau/8 = 8$--$32$\,B, and V1's stored Merkle tree a further \akmTreeBytes{}\,B.
For V2 at $\tau=128$ the total is \akmVtwoPctSmall{}\% of a 200\,B payload and \akmVtwoPctLarge{}\% of a \akmLargeSize{}\,B payload.
The commitment that buys detection of wrong-key faults is the cheapest component of the audit; P7 held.

\begin{table}[t]
  \centering
  \caption{Experiment 1: detection rate per fault and verifier (100 seeds per cell; F0 row shows false positives).}
  \label{tab:e1}
  \begin{tabular}{lrrrrr}
    \toprule
    Fault & V0 & V1 & V2$^{-}$ & V2 & V3 \\
    \midrule
    F0 & FP 0/100 & FP 0/100 & FP 0/100 & FP 0/100 & FP 0/100 \\
    F1a & 100\% & 100\% & 100\% & 100\% & 100\% \\
    F1b & 100\% & 100\% & 100\% & 100\% & 100\% \\
    F2 & 100\% & 100\% & 100\% & 100\% & 100\% \\
    F3 & 100\% & 100\% & 100\% & 100\% & 100\% \\
    F4r & 100\% & 100\% & 100\% & 100\% & 100\% \\
    F4m & 100\% & 100\% & 100\% & 100\% & 100\% \\
    F5 & 0\% & 0\% & 100\% & 100\% & 100\% \\
    F5 @ $\varepsilon$=0.05 & 0\% & 0\% & 100\% & 100\% & 100\% \\
    F5b & 0\% & 0\% & 100\% & 100\% & 100\% \\
    F5c & 0\% & 0\% & 100\% & 100\% & 100\% \\
    F6 & 100\% & 100\% & 100\% & 100\% & 100\% \\
    F7 & n/a & n/a & 0\% & 100\% & n/a \\
    \bottomrule
  \end{tabular}
\end{table}

\begin{table}[t]
  \centering
  \caption{RQ4: V1-S detection at $k=100$ (sampling mode). Theory: $1-(1-m/N)^k$.}
  \label{tab:rq4}
  \begin{tabular}{llrrrrr}
    \toprule
    Fault & $\varepsilon$ & $m$ & detected & 95\% CI & theory & trials \\
    \midrule
    F1b & 0.1\% & 10 & 0.091 & [0.075, 0.110] & 0.095 & 1000 \\
    F1b & 1\% & 100 & 0.634 & [0.604, 0.663] & 0.634 & 1000 \\
    F1b & 5\% & 500 & 0.995 & [0.988, 0.998] & 0.994 & 1000 \\
    F3 & 0.1\% & 10 & 0.103 & [0.086, 0.123] & 0.095 & 1000 \\
    F3 & 1\% & 100 & 0.616 & [0.585, 0.646] & 0.634 & 1000 \\
    F3 & 5\% & 500 & 0.994 & [0.987, 0.997] & 0.994 & 1000 \\
    F5 & 0.1\% & 10 & 0.000 & [0.000, 0.004] & 0 (Cor. 1.1) & 1000 \\
    F5 & 1\% & 100 & 0.000 & [0.000, 0.004] & 0 (Cor. 1.1) & 1000 \\
    F5 & 5\% & 500 & 0.000 & [0.000, 0.004] & 0 (Cor. 1.1) & 1000 \\
    \bottomrule
  \end{tabular}
\end{table}

\begin{table}[t]
  \centering
  \caption{Experiment 2: migration wall-clock seconds, median (IQR) over the measured reps.}
  \label{tab:e2time}
  \begin{tabular}{rrrrrr}
    \toprule
    $N$ & B & no proof, rec. & proof, rec. & no proof, group & proof, group \\
    \midrule
    1{,}000 & 200 & 20.176 (0.249) & 25.245 (2.082) & 0.264 (0.007) & 0.296 (0.005) \\
    1{,}000 & 2048 & 22.916 (2.605) & 34.218 (1.243) & 0.277 (0.007) & 0.319 (0.008) \\
    10{,}000 & 200 & 266.141 (5.571) & 345.348 (3.224) & 2.969 (0.030) & 3.106 (0.039) \\
    10{,}000 & 2048 & 287.831 (0.271) & 355.838 (3.652) & 3.029 (0.068) & 3.135 (0.064) \\
    \bottomrule
  \end{tabular}
\end{table}

\begin{table}[t]
  \centering
  \caption{Experiment 2: verification wall-clock seconds, median (IQR).}
  \label{tab:e2verify}
  \begin{tabular}{rrrrrr}
    \toprule
    $N$ & B & V0 & V1-S & V2$^{-}$ & V2 \\
    \midrule
    1{,}000 & 200 & 0.058 (0.002) & 0.020 (0.001) & 0.150 (0.003) & 0.106 (0.001) \\
    1{,}000 & 2048 & 0.067 (0.001) & 0.021 (0.003) & 0.151 (0.003) & 0.113 (0.001) \\
    10{,}000 & 200 & 0.544 (0.004) & 0.062 (0.001) & 1.367 (0.011) & 0.967 (0.004) \\
    10{,}000 & 2048 & 0.587 (0.019) & 0.058 (0.002) & 1.390 (0.017) & 1.019 (0.016) \\
    \bottomrule
  \end{tabular}
\end{table}

\begin{table}[t]
  \centering
  \caption{P4: V2 verification time divided by proof-emitting migration time (pre-registered threshold 0.10; $^{\ast}$ = below 0.10).}
  \label{tab:p4}
  \begin{tabular}{rrrrrr}
    \toprule
    $N$ & B & rec.\ wall & rec.\ CPU & group wall & group CPU \\
    \midrule
    1{,}000 & 200 & 0.004$^{\ast}$ & 0.067$^{\ast}$ & 0.357 & 0.434 \\
    1{,}000 & 2048 & 0.003$^{\ast}$ & 0.070$^{\ast}$ & 0.354 & 0.431 \\
    10{,}000 & 200 & 0.003$^{\ast}$ & 0.061$^{\ast}$ & 0.311 & 0.389 \\
    10{,}000 & 2048 & 0.003$^{\ast}$ & 0.063$^{\ast}$ & 0.325 & 0.404 \\
    \bottomrule
  \end{tabular}
\end{table}

\begin{table}[t]
  \centering
  \caption{Storage overhead per record (bytes), from real counters.}
  \label{tab:storage}
  \begin{tabular}{rrlrrrrr}
    \toprule
    B & $\tau$ & verifier & commit. & receipts & tree & total & \% payload \\
    \midrule
    200 & 64 & V2 & 8.0 & 83.0 & 0.0 & 91.0 & 45.50\% \\
    200 & 96 & V2 & 12.0 & 83.0 & 0.0 & 95.0 & 47.50\% \\
    200 & 128 & V0 & 0.0 & 83.0 & 0.0 & 83.0 & 41.50\% \\
    200 & 128 & V1 & 0.0 & 83.0 & 32.0 & 115.0 & 57.51\% \\
    200 & 128 & V2$^{-}$ & 0.0 & 83.0 & 0.0 & 83.0 & 41.50\% \\
    200 & 128 & V2 & 16.0 & 83.0 & 0.0 & 99.0 & 49.50\% \\
    200 & 256 & V2 & 32.0 & 83.0 & 0.0 & 115.0 & 57.50\% \\
    2048 & 64 & V2 & 8.0 & 83.0 & 0.0 & 91.0 & 4.44\% \\
    2048 & 96 & V2 & 12.0 & 83.0 & 0.0 & 95.0 & 4.64\% \\
    2048 & 128 & V0 & 0.0 & 83.0 & 0.0 & 83.0 & 4.05\% \\
    2048 & 128 & V1 & 0.0 & 83.0 & 32.0 & 115.0 & 5.62\% \\
    2048 & 128 & V2$^{-}$ & 0.0 & 83.0 & 0.0 & 83.0 & 4.05\% \\
    2048 & 128 & V2 & 16.0 & 83.0 & 0.0 & 99.0 & 4.83\% \\
    2048 & 256 & V2 & 32.0 & 83.0 & 0.0 & 115.0 & 5.62\% \\
    \bottomrule
  \end{tabular}
\end{table}

